% ============================================================================
% Fichier  : partie1/P01_C01_forensique_numerique.tex : FORN (LPI342)
% Sources : ISO/IEC 27037, 27042, 27043 ; Locard ; ACPO. Opération TSANAGA.
% ============================================================================

\chapter{Qu'est-ce que la Forensique Numérique}
\label{chap:forensique}
\chaptermark{La forensique numérique}

\epigraph{%
    « La trace numérique est fragile~; c'est le soin de l'enquêteur\\
    qui la rend solide. »%
}{Principe du cours}

\niveauChapitre{Fondamental}{Le prologue~; notions de systèmes et de réseaux.}

\begin{boxobjectifs}
\begin{itemize}
    \item Définir la forensique numérique et ses termes voisins.
    \item Situer son histoire et ses figures fondatrices.
    \item Distinguer ses domaines (disque, mémoire, réseau, mobile, nuage).
    \item Comprendre son double ancrage, scientifique et juridique.
    \item Situer les référentiels (ISO/IEC 27037, 27042, 27043, ACPO).
\end{itemize}
\end{boxobjectifs}

\bigskip

% ============================================================================
\section*{Question directrice : qu'est-ce qui distingue l'enquêteur numérique d'un informaticien ordinaire ?}
% ============================================================================

Un informaticien sait lire un disque~; un enquêteur forensique le lit
\emph{sans l'altérer}, documente chaque geste, et sépare ce qu'il constate de
ce qu'il en déduit. Ce chapitre définit la discipline, son histoire et ses
domaines.

\medskip
\textbf{Plan du chapitre.}
\begin{enumerate}
    \item Définir la discipline (section~\ref{sec:c01-definir}).
    \item Histoire et figures (section~\ref{sec:c01-histoire}).
    \item Les domaines (section~\ref{sec:c01-domaines}).
    \item Le double ancrage (section~\ref{sec:c01-ancrage}).
\end{enumerate}

\separateur

% ============================================================================
\section{Définir la Discipline}
\label{sec:c01-definir}
% ============================================================================

\begin{boxdef}
\termecle{Investigation numérique}~: démarche générale de collecte, d'analyse
et d'interprétation d'éléments numériques pour éclairer une situation, dans
tout contexte (judiciaire, administratif, privé). La \termecle{forensique
numérique} (ou criminalistique numérique) en est le sous-ensemble à vocation
probatoire, soumis à la rigueur de la preuve.
\end{boxdef}

\begin{tableaucours}{L{4.4cm}L{7.8cm}}{%
    \tableauHeader{Terme} &
    \tableauHeader{Sens}
}
Investigation numérique & Démarche générale, tous contextes \\
Forensique (digital forensics) & Cadre probatoire, rigueur de la preuve \\
Criminalistique numérique & Insiste sur la dimension scientifique \\
Recherche de la preuve & Finalité juridique \\
\end{tableaucours}
\vspace{-0.8em}
\begin{center}{\small\itshape Tableau~1.1. Des termes voisins}\end{center}

% ============================================================================
\section{Histoire et Figures}
\label{sec:c01-histoire}
% ============================================================================

La forensique numérique hérite des sciences de la trace. \termecle{Edmond
Locard} (France, 1910) en énonce le principe fondateur~: « tout contact
laisse une trace »~\cite{locard1920}. Appliqué au numérique, il signifie que
toute action sur un système y laisse des marques, pour qui sait les lire et
les préserver.

% ============================================================================
\section{Les Domaines}
\label{sec:c01-domaines}
% ============================================================================

\begin{tableaucours}{L{3.4cm}L{8.8cm}}{%
    \tableauHeader{Domaine} &
    \tableauHeader{Objet}
}
Disque & Systèmes de fichiers, fichiers supprimés, artefacts \\
Mémoire & État vivant du système, secrets, code \\
Réseau & Flux, exfiltration, pilotage, chronologie \\
Mobile & Téléphones~: messages, applications, géolocalisation \\
Nuage et applications & Données hébergées, journaux de services \\
\end{tableaucours}
\vspace{-0.8em}
\begin{center}{\small\itshape Tableau~1.2. Les domaines de la forensique numérique}\end{center}

% ============================================================================
\section{Le Double Ancrage}
\label{sec:c01-ancrage}
% ============================================================================

La forensique numérique s'appuie sur des référentiels qui lui donnent sa
rigueur~: l'\norme{ISO/IEC 27037} (acquisition et préservation),
l'\norme{ISO/IEC 27042} (analyse et interprétation), l'\norme{ISO/IEC 27043}
(principes de l'investigation), et le guide \sigle{ACPO}. Elle s'ancre aussi
dans le droit (Partie~II).

\begin{boxscenario}{Opération TSANAGA. Une discipline, pas une intuition}
\textbf{Constatation.} Devant l'ordinateur saisi, M.~BIKORO ne l'allume pas~:
il le photographie scellé, note son état, et prépare l'acquisition selon
l'ISO/IEC 27037.

\medskip
\textbf{Analyse} (confiance élevée). Le réflexe « ouvrir pour voir » aurait
altéré la preuve. La discipline impose la méthode avant la curiosité~: c'est
ce qui sépare l'enquête de la manipulation.
\end{boxscenario}

\begin{boxmetier}{Analyste forensique}
\textbf{La méthode avant l'outil.} L'analyste connaît les outils, mais c'est
sa méthode (préserver, documenter, séparer fait et interprétation) qui donne
à son travail une valeur probante.
\end{boxmetier}

\separateur

\begin{boxresume}
\textbf{Ce qu'il faut retenir de ce chapitre~:}
\begin{itemize}
    \item \textbf{La forensique} est l'investigation numérique à vocation
          probatoire.
    \item \textbf{Locard}~: « tout contact laisse une trace ».
    \item \textbf{Cinq domaines}~: disque, mémoire, réseau, mobile, nuage.
    \item \textbf{Double ancrage}~: scientifique (ISO 27037/42/43, ACPO) et
          juridique.
    \item \textbf{La méthode} distingue l'enquêteur de l'informaticien
          ordinaire.
\end{itemize}
\end{boxresume}

\begin{boxinfo}
\textbf{Points de contrôle de l'élève.} Je sais~: définir la forensique et
ses termes voisins~; citer Locard~; nommer les domaines~; situer les
référentiels~; dire ce qui distingue l'enquêteur.
\end{boxinfo}

\bigskip

\begin{boxexo}[title={\ding{45}\quad Exercice~1.1. Termes \niveaubase}]
Distinguez investigation numérique et forensique numérique.
\end{boxexo}

\begin{boxexo}[title={\ding{45}\quad Exercice~1.2. Locard \niveaubase}]
Énoncez le principe de Locard et son sens pour le numérique.
\end{boxexo}

\begin{boxexo}[title={\ding{45}\quad Exercice~1.3. Domaines \niveaubase}]
Citez les cinq domaines de la forensique numérique.
\end{boxexo}

\begin{boxexo}[title={\ding{45}\quad Exercice~1.4. Pirate ou enquêteur \niveaumoyen}]
Sur un même disque, qu'est-ce qui distingue le geste du pirate de celui de
l'enquêteur~?
\end{boxexo}

\begin{boxexo}[title={\ding{45}\quad Exercice~1.5. Référentiels \niveaumoyen}]
Associez à chaque norme son objet~: 27037, 27042, 27043.
\end{boxexo}

\begin{boxexo}[title={\ding{45}\quad Exercice~1.6. TSANAGA \niveaumoyen}]
Pourquoi M.~BIKORO n'allume-t-il pas l'ordinateur saisi~?
\end{boxexo}

\begin{boxexo}[title={\ding{45}\quad Exercice~1.7. Double ancrage \niveaumoyen}]
Expliquez le double ancrage, scientifique et juridique, de la discipline.
\end{boxexo}

\begin{boxexo}[title={\ding{45}\quad Exercice~1.8. Criminalistique \niveaumoyen}]
Pourquoi parle-t-on de « criminalistique » numérique~?
\end{boxexo}

\begin{boxexo}[title={\ding{45}\quad Exercice~1.9. Posture \niveauexpert}]
Rédigez une courte charte de la posture de l'enquêteur numérique.
\end{boxexo}

\begin{boxexo}[title={\ding{45}\quad Exercice~1.10. Panorama \niveauexpert}]
Pour l'opération TSANAGA, quels domaines seront probablement mobilisés, et
pourquoi~?
\end{boxexo}

\bigskip

\begin{boxinfo}
\textbf{Pour aller plus loin.}
\begin{itemize}
    \item ISO/IEC 27037~\cite{iso27037}, 27043~\cite{iso27043}~;
          ACPO~\cite{acpo2012}~; Locard~\cite{locard1920}.
\end{itemize}
\end{boxinfo}

\begin{boxinfo}
\textbf{Transition.} Toute enquête suit une démarche~: le
Chapitre~\ref{chap:demarche} en présente les phases.
\end{boxinfo}
